\documentclass[11pt]{article}
\usepackage[margin=1in]{geometry}
\usepackage{amsmath,amssymb,amsthm,hyperref}
\newtheorem{theorem}{Theorem}
\newtheorem{lemma}{Lemma}
\title{The polar-pentagon perimeter centroid has degree at least four}
\author{Principia Math In Progress}
\date{21 September 2026}
\begin{document}
\maketitle
\begin{abstract}
We give an exact finite certificate excluding every cubic from the
perimeter-centroid locus of the polar pentagons in the
focus-circle Poncelet family. Algebraicity follows from the tangent
recurrence and the edge-length square roots. Together with the
previous general polar-centroid theorem, this resolves the intended
pentagon conjecture of Bellio, Garcia and Reznik.
\end{abstract}

\section*{Task 1: the original problem and its intended scope}
Consider the convex Poncelet pentagons inscribed in a parabola of
focal distance $f>0$ and circumscribed about a circle centered at its
focus. Their polar polygons have as sides the tangents to the
parabola at the five original vertices. The perimeter centroid is
the centroid of their five segments with positive Euclidean
length as mass, including overlapping segments with multiplicity.

Conjecture 3 of \cite{BGR} asserts straight vertex- and area-centroid
loci and perimeter-centroid degree at least four. Its word
\emph{perpendicular} before \emph{directrix} is a direction error:
the intended direction is parallel, as stated in its general
Conjecture 5. We prove the substantive degree assertion, rather
than count that wording error as a counterexample.

\begin{theorem}
The perimeter-centroid locus is an algebraic locus of degree at
least four: it is contained in an algebraic plane curve, but no
nonzero polynomial of total degree at most three vanishes on it.
The vertex- and area-centroid loci are complete lines parallel
to the directrix.
\end{theorem}
Here degree means the degree of the algebraic closure of the
physical locus. The positive length convention selects real
branches; as in the source's quadrilateral example, an implicit
equation can have additional real points. We claim the lower bound,
not an exact degree or an equation with no extraneous branches.

The line assertions are the $N=5$ case of the previously established
general theorem \cite{polar}; they receive no additional credit.
In coordinates $y^2+4fx=0$, they are $x=h_0$ and
$x=(8f^2-r^2)/(12f)$, respectively. The new proof below establishes
the independent perimeter-degree assertion.

\section*{Algebraic coordinates and the six samples}
A similarity permits $f=1$. By \cite[Proposition 25]{BGR}, the
required radius is the unique positive root $r$ of
\begin{equation}\label{radius}
 r^6+12r^5-28r^4+32r^3+112r^2-64r-64=0,
 \qquad \frac{995}{1000}<r<\frac{996}{1000}.
\end{equation}
It is approximately $0.995219078669842$. Write
$P(t)=(-t^2/4,t)$. The chord through $P(t),P(s)$ is tangent to the
focus circle if and only if
\begin{equation}\label{tangent}
 (ts+4)^2=r^2\bigl(16+(t+s)^2\bigr).
\end{equation}
Indeed, the chord is $4x+(t+s)y-ts=0$ and the focus is $(-1,0)$.

For a starting value $a$, put $t_0=a$ and define
\begin{align}
 A&=a^2-r^2,& B&=(8-2r^2)a,& C&=16-16r^2-r^2a^2,\nonumber\\
 t_1&=\frac{-B+\sqrt{B^2-4AC}}{2A},\label{start}\\
 t_{j+1}&=-\frac{(8-2r^2)t_j}{t_j^2-r^2}-t_{j-1},
       &&j=1,2,3.\label{step}
\end{align}
The positive square root in \eqref{start} chooses one traversal
direction. Equation \eqref{step} is Vieta's formula for the other
root of \eqref{tangent}. The cited closure theorem makes the
resulting five vertices a closed Poncelet pentagon.

The intersection of the tangents at parameters $t_i,t_{i+1}$ is
\begin{equation}\label{polar}
 Q_i=\left(-\frac{t_it_{i+1}}4,\frac{t_i+t_{i+1}}2\right),
 \qquad i\pmod5 .
\end{equation}
Consequently, with positive square roots throughout,
\begin{equation}\label{center}
 \ell_i=\sqrt{(Q_{i+1,x}-Q_{i,x})^2+
                    (Q_{i+1,y}-Q_{i,y})^2},\qquad
 (X,Y)=\frac{\sum_{i=0}^4\ell_i(Q_i+Q_{i+1})}
                    {2\sum_{i=0}^4\ell_i}.
\end{equation}
These formulas show algebraicity. The two possible values of
$t_1$ form a finite algebraic extension of $\mathbb R(a)$;
the other parameters and polar vertices are rational in it.
Adjoining the five square roots in \eqref{center} is again a
finite algebraic extension. Its transcendence degree is one, so
$X,Y$ satisfy a nonzero polynomial relation. There are only finitely
many branches, so multiplying their relations gives one relation
for the entire locus. Exceptional parameter values are covered by
continuity whenever the polygon and its centroid are finite.

Use the following six \emph{exact rational} starting values:
\begin{equation}\label{samples}
 a_1=\frac3{25},\quad a_2=\frac{27}{100},\quad
 a_3=\frac{43}{100},\quad a_4=\frac{17}{25},\quad
 a_5=\frac{131}{100},\quad a_6=\frac{12}{5}.
\end{equation}
Let $(X_i,Y_i)$ be defined exactly by
\eqref{radius}--\eqref{center}. For orientation only, their
approximate coordinates are
\[
\begin{array}{c|r|r}
i&X_i&Y_i\\ \hline
1&0.531703184427& 9.271732856495\\
2&0.501537653846&-1.067231332343\\
3&0.496287921973& 0.232794002460\\
4&0.509767562441& 2.074695123461\\
5&0.517049614999&-3.248106062772\\
6&0.496610704943&-0.328686222928
\end{array}
\]
The rounded entries in this table are not used in the certificate.

\section*{An exact determinant certificate}
Form the matrix with rows
\[
 M_i=(1,\ X_i,\ X_i^2,\ X_i^3,\ Y_i^2,\ X_iY_i^2).
\]
\begin{lemma}\label{certificate}
All six sample polygons are finite, strictly convex members of the
required family; all $Y_i$ are nonzero; and
\begin{equation}\label{detbound}
 -\frac{1769}{10^{18}}<\det M<-\frac{1768}{10^{18}}.
\end{equation}
\end{lemma}
\begin{proof}
The accompanying \texttt{verify.py} is an exact certificate evaluator
using only integers and rational numbers from the Python standard
library. The algorithm and its bounds are specified here so that
the computation does not depend on floating-point behavior.

Set $S=2^{512}$ and represent $[L/S,U/S]$ by the integers $(L,U)$.
For addition and negation use the exact endpoint formulas.
For multiplication or division, take the minimum and maximum
of all four endpoint products or quotients and round outwards
to multiples of $1/S$. Every division asserts that its denominator
interval excludes zero. For a nonnegative interval, the square-root
endpoints in integer units are
\[
 \left\lfloor\sqrt{LS}\right\rfloor,\qquad
 \left\lceil\sqrt{US}\right\rceil .
\]
Integer square root and an exact square comparison implement them.
Squaring uses zero as the lower endpoint when the input crosses zero.
These rules preserve enclosure at every operation.

The polynomial in \eqref{radius} has opposite signs at $995/1000$
and $996/1000$. Evaluating its derivative by interval arithmetic on
this whole interval gives a strictly positive lower bound, so
the root there is unique. Exact rational bisection refines the
root bracket to width below $2^{-520}$; its outward dyadic
enclosure is used in every sample.

Evaluate \eqref{start}--\eqref{center}, asserting nonzero
denominators, positive radicands, distinct vertex parameters and
$Y_i\ne0$. Exact signs of edge cross products show that every
other original vertex and the focus lie strictly on the same
side of each directed edge. Thus the original pentagons are
strictly convex with the circle on their interior side.
Closure uses Proposition 25, not a small numerical residual.
All polar segment lengths are certified positive.

Finally, evaluate the determinant as the signed sum of all $6!=720$
permutation products with these same interval rules. The resulting
enclosure, rounded outwards here to 35 decimal places, is
\[
 \bigl[-0.00000000000000176874062193541107867,\,
       -0.00000000000000176874062193541107866\bigr].
\]
It lies strictly within \eqref{detbound}. The full enclosure,
sample bounds and all assertions are recorded in
\texttt{verification-results.json}. This proves the lemma by a
finite exact calculation.
\end{proof}

\section*{Why the certificate excludes every cubic}
Reflection in the parabola's axis preserves the Poncelet family,
polarity, and positive Euclidean edge lengths. Hence the locus
contains both $(X_i,Y_i)$ and $(X_i,-Y_i)$.

Suppose a polynomial $H(X,Y)$ of total degree at most three
vanishes on the entire locus. Its even part in $Y$ has the form
\[
 H_{\rm even}=c_0+c_1X+c_2X^2+c_3X^3+c_4Y^2+c_5XY^2 .
\]
Adding its values at the reflected samples gives $Mc=0$.
Lemma \ref{certificate} makes $M$ invertible, so all six
coefficients vanish. The remaining odd part is
\[
 H_{\rm odd}=Y(d_0+d_1X+d_2X^2+d_3Y^2).
\]
Since each $Y_i\ne0$, the six sample equations say that a linear
combination of columns $1,2,3,5$ of $M$ is zero. Those columns
are independent, again because $M$ is invertible. Thus all four
$d_j$ vanish, and $H=0$.

The algebraic locus therefore has degree at least four.
This also proves it is nonconstant. Restoring $f$ applies an
invertible similarity, which preserves degree and completes the
perimeter assertion for every scale.

\section*{Four proposed follow-up tasks}
\begin{enumerate}
 \item \textbf{Task 2.} Determine the exact degree and a minimal
 implicit equation of this perimeter-centroid locus.
 \item \textbf{Task 3.} Determine its real branches and which
 parts are attained by finite polar pentagons.
 \item \textbf{Task 4.} Find exact bounds for its horizontal
 coordinate and identify all extremizing configurations.
 \item \textbf{Task 5.} Determine the degree of the polar
 perimeter-centroid locus for each convex period $N\geq6$.
\end{enumerate}

\begin{thebibliography}{9}
\bibitem{BGR} F. Bellio, R. Garcia, D. Reznik,
\emph{Parabola-Inscribed Poncelet Polygons Derived from the Bicentric Family},
KoG \textbf{26} (2022), 3--15,
\href{https://doi.org/10.31896/k.26.1}{doi:10.31896/k.26.1};
\href{https://arxiv.org/abs/2111.00979}{arXiv:2111.00979v3}.
\bibitem{polar} Principia Math In Progress,
\emph{Straight centroid loci for polar Poncelet polygons},
MathDB 352456, 21 September 2026.
Repository directory \texttt{problems/352456}, with the full proof
in \texttt{solution.tex} and \texttt{solution.pdf},
commit \texttt{ece807501ca0}.
\end{thebibliography}
\end{document}
